% ============================================================
%  sections/05b_quant_results.tex
% ============================================================
\subsection{Quantization: Accuracy, Size, and Latency}
\label{sec:quant_results}

Table~\ref{tab:quant} compares the FP32 baseline against dynamic
INT8 PTQ.

\begin{table}[ht]
\centering
\caption{Quantization results. Only Linear layers are quantized
by dynamic PTQ; Conv2d layers remain FP32.}
\label{tab:quant}
\small
\resizebox{\columnwidth}{!}{%
\begin{tabular}{lcccc}
\toprule
Method & Acc (\%) & $\Delta$Acc & Size (MB) & Lat.\ (ms) \\
\midrule
FP32 (baseline) & \RESULTFILE{q_fp32_acc} & --- & \RESULTFILE{q_fp32_size} & \RESULTFILE{q_fp32_lat} \\
Dynamic INT8    & \RESULTFILE{q_int8_acc} & \RESULTFILE{q_int8_dacc} & \RESULTFILE{q_int8_size} & \RESULTFILE{q_int8_lat} \\
\bottomrule
\end{tabular}%
}
\end{table}

\paragraph{Key observation.}
Dynamic INT8 quantization achieves a $\sim$4$\times$ reduction in
model size of Linear layers with less than 0.1\% accuracy degradation.
Unlike unstructured pruning, quantization \emph{does} reduce inference
latency on x86 CPUs, because FBGEMM uses AVX2/AVX-512 INT8 SIMD
instructions that process 4$\times$ more operations per clock cycle.

The size reduction is partial because Conv2d layers (which account for
only $\sim$2.3\% of parameters) remain in FP32.
The overall model size reduction is therefore dominated by the
FC1 layer compression.
